\section{Célközönség és perszónák}

A célközönség meghatározása a termékkoncepció egyik fontos része, mert a fejlesztési döntések nem választhatók el attól, hogy kiknek készül az alkalmazás. A PRD alapján a termék elsődleges célcsoportját azok az alkalmi filmnézők alkotják, akik havonta néhány filmet néznek, és egyszerű, gyors megoldást keresnek a filmjeik nyilvántartására. Másodlagos célcsoportként megjelennek azok a filmrajongók is, akik tudatosabban követik filmnézési szokásaikat, és hosszabb távon részletesebb statisztikákat vagy közösségi funkciókat is használnának.

Az első célszemély az alkalmi filmkövető típusa. Ő rendszeresen néz filmeket streaming szolgáltatásokon vagy moziban, de nem feltétlenül tekinti magát filmrajongónak. Számára a legfontosabb érték az egyszerűség, gyorsan szeretné rögzíteni, hogy mit látott már, mit szeretne később megnézni, és esetleg milyen benyomása volt egy adott filmről. Ez a felhasználói típus indokolja, hogy az MVP ne legyen túlterhelt bonyolult funkciókkal, hosszú értékelési lehetőségekkel vagy összetett közösségi funkciókkal. A keresésnek, a megnézendő filmek listájának kezelésének és a megnézett filmek naplózásának ezért kevés lépésből kell állnia.

A második célszemély a közösségi filmnéző, aki nemcsak saját magának vezetne listát, hanem kíváncsi arra is, hogy ismerősei milyen filmeket néznek és hogyan értékelik azokat. Ennél a felhasználói típusnál a filmnézés egy bizonyos szociális aspektussá is válik. A PRD-ben szereplő későbbi közösségi funkciók, például a publikus profil, a követési rendszer és az aktivitási feed ehhez a célcsoporthoz kapcsolódik. Ezek a funkciók nem feltétlenül részei az első MVP-nek, de fontosak a termék számára hosszútávon.

A harmadik célszemély a tudatos filmrajongó vagy filmkedvelő, aki nagyobb mennyiségben néz filmeket, szívesen értékel és rendszerez, de nem feltétlenül szeretne hosszú kritikákat írni. Számára a személyes profil, a megtekintett filmek listája, az értékelések és a későbbi statisztikai kimutatások bizonyulhatnak értékesnek. Ez a perszóna bemutatja, hogy a termékvízióban a filmnapló mellett megjelenhetnek mélyebb funkciók is, például műfaji statisztikák, éves összegzések, egyéni listák vagy ajánlási lehetőségek.

A célszemélyek közös jellemzője, hogy mindhárom esetben fontos a gyors használhatóság és az átlátható felület. A különbség inkább abban jelenik meg, hogy az egyes felhasználói típusok milyen mélységben használnák az alkalmazást. Az alkalmi felhasználó számára az MVP alapfunkciói is elegendő értéket adhatnak, míg a közösségi filmnéző és a filmrajongó inkább a későbbi verziókban megjelenő bővítéseket tartaná hasznosnak. Ez a felosztás segít abban, hogy az első fejlesztési szakasz a legszélesebb célcsoport számára hasznos funkciókra koncentráljon, miközben a termék jövőbeli bővíthetőségére is lehetőség adódik.

Ennek alapján a WatchDB célközönsége nem egyetlen, szűken meghatározott, hanem több, eltérő igényű és filmnézési szokásokkal rendelkező csoportból áll. A PRD-ben bemutatott célszemélyek szerepe ezért az, hogy a funkciók fontossági sorrendje először a legnagyobb célközönség alapvető igényeihez igazodjon, és csak ezt követően kerüljenek előtérbe a szűkebb felhasználói csoportokat kiszolgáló bővítések.
